\chapter{پیشینه و مفاهیم لازم}
\label{ch:background}

این فصل تنها مفاهیمی را می‌آورد که برای خواندن روش و نتایج ضروری است، و
جایگاه پروژه را در برابر کارهای پیشین مشخص می‌کند.

\section{دو خانواده معماری}

سامانه‌های گفت‌وگوی گفتاری در دو خانواده قرار می‌گیرند. در خانواده
\مهم{آبشاری}، بازشناسی گفتار، موتور گفت‌وگوی متنی و تبدیل متن به گفتار
جداگانه اجرا می‌شوند. این معماری قابل ممیزی و قابل جایگزینی است، اما تأخیر
مؤلفه‌ها جمع می‌شود و خطای هر مرحله به مرحله بعد منتقل می‌گردد. در خانواده
\مهم{سرتاسری}، مدل مستقیماً از گفتار به گفتار می‌رود. مدل‌های تمام‌دوطرفه
ذاتی، مانند \en{Moshi}، دو جریان صوتی را به‌صورت موازی مدل می‌کنند و برای
نوبت‌گیری به بخش‌بندی صریح گفتار نیاز ندارند.

تعریف این پروژه روح سرتاسری دارد، اما بودجه کارشناسی و نبود آموزش‌دهنده
رسمی فارسی، مسیر آبشاری را به‌عنوان سامانه تحویل‌شده اجتناب‌ناپذیر کرد.

\section{مدل زبانی، تطبیق کم‌پارامتر، و اریبی مواجهه}

پاسخ‌گوی متنی این پروژه یک مدل دستور‌پیرو است. تنظیم دقیق کامل یک مدل
چند‌میلیارد‌پارامتری روی سخت‌افزار محدود عملی نبود؛ بنابراین مسیر مستقیم از
\en{LoRA}~\cite{hu2022lora} استفاده کرد: به‌روزرسانی وزن به‌صورت حاصل دو
ماتریس رتبه پایین روی وزن منجمد اضافه می‌شود.

نکته تعیین‌کننده برای تفسیر نتایج مسیر مستقیم، تفاوت آموزش و استنتاج است.
در آموزش، مدل در هر گام نشانه مرجع پیشین را می‌بیند؛ این شیوه \مهم{تحمیل
معلم} است و زیان گزارش‌شده تحت همین شرط محاسبه می‌شود. در استنتاج، مدل
باید بر نشانه‌های تولیدی خودش شرط بگذارد. اگر این دو توزیع فاصله بگیرند،
خطا تجمعی می‌شود. این ناسازگاری \مهم{اریبی مواجهه} نام دارد~\cite{ranzato2016scheduledsampling}
و راهکار متعارف آن نمونه‌گیری
زمان‌بندی‌شده است~\cite{bengio2015scheduledsampling}. پیامد عملی صریح است:
کاهش زیان تحت تحمیل معلم، تولید خودرگرسیو معتبر را تضمین نمی‌کند.

\section{بازشناسی، گفتارسازی، و ویژگی‌های وقفه}

بازشناسی گفتار سامانه، یک مدل ترکیبی \en{RNNT-CTC} از خانواده
\en{Conformer} است که به‌صورت محلی با پشته
\en{NeMo}~\cite{kuchaiev2019nemo} اجرا می‌شود. تبدیل متن به گفتار با موتور
\en{Piper}~\cite{piper2023} و مدل صوتی فارسی انجام می‌گیرد.

آشکارساز وقفه ویژگی‌های کلاسیک را به‌صورت غلتان محاسبه می‌کند: انرژی، نرخ
گذر از صفر، ضرایب \en{MFCC}~\cite{davis1980mfcc} و بسامد پایه با روش
\en{YIN}~\cite{decheveigne2002yin}. طبقه‌بند، تقویت گرادیانی
است~\cite{friedman2001gbm}. فاصله‌های اطمینان دقت با روش
ویلسون~\cite{wilson1927interval} و در سطح نشست به‌صورت بلوکی گزارش
می‌شوند. در گفت‌وگوی انسانی، شکاف‌های نوبت‌گیری اغلب در حدود
\N{200} میلی‌ثانیه است~\cite{skantze2021turntaking}؛ بنابراین تصمیم وقفه
باید بسیار سریع‌تر از تولید کامل پاسخ گرفته شود.

\section{جایگاه در کارهای پیشین}
\label{sec:positioning}

\en{Mini-Omni2} و \en{LLaMA-Omni2} مخزن استنتاج دارند اما آموزش‌دهنده کامل
و رسمی گفتار‌به‌گفتار ارائه نمی‌دهند. \en{Freeze-Omni} تمام‌دوطرفه است اما
آموزش‌دهنده آن برای زبان جدید کامل نیست.
\en{Qwen2.5-Omni} فارسی را در فهرست ورودی و خروجی گفتاری ندارد و نسخه
هفت‌میلیارد‌پارامتری آن از بازه حافظه هدف خارج است.
\en{PersonaPlex}~\cite{roy2026personaplex} بر پایه \en{Moshi} نقش و صدا را
شرطی می‌کند، اما در زمان تصمیم این پروژه آموزش‌دهنده رسمی و در دسترس برای
زبان جدید فراهم نبود. در میان مدل‌های بررسی‌شده، تنها خانواده \en{Moshi}
هر سه شرط مسیر مستقیم را داشت: تمام‌دوطرفه ذاتی، وزن متن‌باز، و
آموزش‌دهنده رسمی تطبیق کم‌پارامتر.

برای فارسی، منابع خوانشی تک‌گوینده جایگزین پیکره مکالمه‌ای تعاملی نیستند.
جدول~\ref{tab:positioning} این مقایسه را خلاصه می‌کند.

\begin{table}[ht]
\centering
\caption{مقایسه این پروژه با کارهای پیشین از نظر قابلیت‌های تعریف پروژه}
\label{tab:positioning}
\small
\begin{tabular}{L{2.6cm}ccccc}
\toprule
\textbf{کار} & \textbf{تمام‌} & \textbf{وزن} & \textbf{آموزش‌دهنده} &
\textbf{پشتیبانی} & \textbf{بازه} \\
             & \textbf{دوطرفه} & \textbf{متن‌باز} & \textbf{رسمی} &
\textbf{فارسی} & \textbf{حافظه} \\
\midrule
\en{Mini-Omni2}   & فرمان‌محور & دارد & ندارد & ندارد & مناسب \\
\en{LLaMA-Omni2}  & ندارد & دارد & ندارد & ندارد & مناسب \\
\en{Freeze-Omni}  & دارد & دارد & جزئی & ندارد & مناسب \\
\en{Qwen2.5-Omni} & دارد & دارد & ندارد & ندارد & نامناسب \\
\en{Moshi}        & ذاتی & دارد & دارد & ندارد & نامناسب \\
\en{PersonaPlex}  & ذاتی & دارد & جزئی & ندارد & نامناسب \\
\midrule
این پروژه، مسیر آبشاری & کنترل بیرونی & دارد & بی‌نیاز & دارد & مناسب \\
این پروژه، مسیر مستقیم & ذاتی & دارد & دارد & ناموفق & نامناسب \\
\bottomrule
\end{tabular}

\vspace{0.5em}
{\footnotesize\itshape ستون «بازه حافظه» تناسب با \N{12} تا \N{24} گیگابایت
تعریف پروژه است، نه کیفیت علمی مدل. «ناموفق» در مسیر مستقیم به معنای شکست
شرط اجرای خودرگرسیو است، نه کاهش‌نیافتن زیان.\par}
\end{table}

\section{سنجه‌ها و داوری مدل}

کیفیت مکانیکی زنجیره با عبور ردیف از شروط رونویسی فارسی، پاسخ فارسی و صوت
غیرساکت سنجیده می‌شود. حفظ محتوای گفتار تولیدشده با بازبازشناسی و نرخ خطای
نویسه و واژه تقریب زده می‌شود؛ این سنجه فهم‌پذیری انسانی نیست. کیفیت معنایی
پاسخ با داوری مدل زبانی روی ارتباط و انسجام، به‌صورت دوسوکور و با دو تکرار،
اندازه گرفته شد. چون داور و نامزد از یک خانواده‌اند، نتیجه در کلاس «تقریب
خودکار هم‌خانواده» می‌ماند و جانشین ارزیابی انسانی یا معیار
\en{ITU-T P.800}~\cite{itu1996p800} نیست. تفکیک آموزش، اعتبارسنجی و آزمون
همیشه بر پایه شناسه نشست منبع است تا نشت گروهی رخ ندهد.
